\chapter{Universal Education}
\label{ch:11}

The third part of the book concerns institutions, and it begins with education because education is the institution through which a society reproduces its capacity for self-direction. The place of education in a program of human direction is not that of one application among others. If the argument of the earlier chapters is accepted, the capacity to direct, which is to say the capacity to understand the options, to contest the claims of those who propose them, and to bear responsibility for choices, is the scarce resource, and a society whose members cannot exercise it has nothing to which capable systems could be subordinate. Later reviews found a smaller effect, of the order of $0.8$ standard deviations both for human tutoring and for well-built intelligent tutoring systems~\cite{vanlehn2011}, so that the original figure is best read as an upper bound. The educational question is therefore twofold. It asks how the instruments described in the second part can extend to every person an education now available to few, and it asks how that extension can be arranged so that it strengthens and does not displace the human relations in which character and judgment are formed. The instructional ideas on which the argument draws, in particular that a tutor supports a learner by taking over only those parts of a task that the learner cannot yet perform, originate in the study of tutoring~\cite{wood1976}.

The empirical foundation for the first aim is the finding, reported by Bloom in 1984~\cite{bloom1984}, that students who received one-to-one tutoring together with mastery-based progression performed on average about two standard deviations above students in conventional classes. The finding was presented as a challenge to educational research, which Bloom called the two sigma problem: to find methods of group instruction as effective as tutoring, since individual tutoring could not be provided to all. Later syntheses have reported smaller effects, of the order of eight-tenths of a standard deviation, for both human tutors and carefully designed computer-based tutors, and the difference between the original and later estimates is a reminder that the effect depends on implementation and measurement. Even the smaller estimate is large by the standards of educational interventions. The argument for the use of capable systems in education begins here: if the benefit of individual attention is substantial and the supply of tutors is limited, then a system that can offer something approximating individual attention at low marginal cost addresses the main limiting factor, and does so in a way that is especially valuable for the children who now receive least.

The case should nonetheless be stated with the qualifications that the evidence warrants. Tutoring effects have been measured in settings where the tutor, human or automated, is part of an organized course of study with defined objectives, assessments that detect mastery, and a framework within which the student is motivated to attend. A conversational system that answers any question on demand is not thereby a tutor in this sense. Without structure it may encourage the substitution of answer-seeking for understanding, a danger recognized in the literature on learning, which distinguishes the performance a student shows while receiving help from the learning that persists when help is withdrawn. A system designed for education should therefore be evaluated by the retention and transfer of what is learned, assessed without assistance after an interval, and not by the satisfaction of students or the correctness of the answers they receive. It should be designed to require effort of the student, to give hints before solutions, and to hold the student responsible for explaining the reasoning, since the evidence supports the view that productive struggle is part of the mechanism by which learning occurs.

The second aim, the preservation of human relations, determines the allocation of tasks between machines and persons, and it is here that the principle of the book bears directly. Some of what happens in a school is instruction, the transmission of knowledge and skill, and much of it can be supervised or delegated. Some of it is assessment of the kind that decides a person's future, and this belongs in the category of judgment that is not delegated. Some of it is the formation of habits, the encouragement of the discouraged, the recognition of a child in difficulty at home, the example of an adult who takes the child seriously, and the initiation into a community of practice, and these are functions that depend on the presence of a person who can be trusted and who is accountable. They are not functions that a machine performs less well and that might one day be performed equally well. They are functions whose value depends on the fact that a person is performing them, in the same way that the value of a promise depends on its being made by someone who can be held to it. The program therefore reserves them to human beings and treats the time that automation frees from instruction as time to be devoted to them.

This leads to a quantitative question on which the program's coherence depends. If instruction is automated to a great extent, how many human mentors does a school require? The answer, derived in the formalization, is that the number is bounded below by the non-delegable portion of human attention per student and does not fall to zero however much instruction is automated. The point is the educational counterpart of a well-known observation in the study of parallel computation, that the speedup of a process is limited by the part that cannot be sped up. It runs against both expectations the public discussion encourages. It contradicts the view that automation of instruction makes teachers unnecessary, and also the view that it will make teaching less demanding. What it predicts is that the role of the teacher will shift toward mentoring, assessment, and the supervision of systems, that the number of teachers required per student will be reduced less than the share of instructional time automated would suggest, and that the training of teachers should change accordingly. It also predicts that the most important investment in an educational transition is in the number and the working conditions of the persons who perform the non-delegable functions, since these constitute the binding constraint.

Questions of governance follow. The curriculum of a system that teaches millions of students is an instrument of extraordinary influence, and its content, emphases, and omissions are choices that ought to be made by processes that the public can contest. The interview notes that systems built from large collections of text reproduce the biases of those texts and the perspectives of the people most heavily represented in them, and the observation applies with force to educational content, where the stakes include the formation of children's picture of the world. The program accordingly requires that educational content be authored or approved through documented procedures with the participation of teachers, parents, and subject specialists, that the sources of generated material be disclosed, and that independent evaluators have access to the system for the purpose of testing it for bias and error. It requires further that the data collected from students be used solely for their education, be subject to deletion on request, and not be used to profile children for commercial purposes. A system that fails these requirements may teach effectively and still be unacceptable, since its effects on the interests of children are mediated by an organization whose own interests are not those of the children.

\section*{Formalization}

Let a school serve $N_s$ students and employ $N_t$ human mentors, each providing $H$ hours of contact per week. For each student, weekly educational activity decomposes into hours of three types: judgment-type activity $h_J$ (assessment, mentoring, safeguarding), which by the principle of the book is non-delegable; supervised activity $h_U$, which can be performed by a system under human oversight; and activity $h_L$ performed by a system without moment-to-moment oversight. Suppose a fraction $a\in[0,1]$ of $h_U$ is performed by the system, the remaining $(1-a)h_U$ by a mentor. The human attention required per student per week is
\[
h(a)=h_J+(1-a)h_U .
\]

\begin{proposition}
The mentor requirement is $N_t\ge N_sh(a)/H$, so the student-to-mentor ratio is at most $H/h(a)$. The function $h(a)$ is decreasing in $a$ and satisfies $h(a)\ge h_J$ for all $a\in[0,1]$, so that no degree of automation of supervised activity raises the student-to-mentor ratio above $H/h_J$. The maximum feasible ratio is $H/h_J$, attained only when $a=1$ or $h_U=0$.
\end{proposition}

\begin{proof}
Total human hours demanded is $N_sh(a)$ and the supply is $N_tH$, whence the inequality. Since $h_U\ge0$ and $a\le1$, $h(a)\ge h_J$ with equality iff $a=1$ or $h_U=0$.
\end{proof}

Numerically, take $H=25$ contact hours, $h_J=1$ and $h_U=4$ hours per week. With no automation $(a=0)$, $h=5$ and the ratio is $5$. With $a=0.75$, $h=2$ and the ratio is $12.5$, so that a country with $10^6$ students requires $80{,}000$ mentors. In the limit $a=1$ the ratio is $25$ and $40{,}000$ mentors are required. Complete automation of supervised work halves the requirement relative to $a=0.75$, but it cannot go lower, and the whole of the saving is bounded by the factor $h(0)/h_J=5$. The proposition thus exhibits the binding role of the non-delegable component.

A second computation concerns the size of the benefit that motivates the effort. If the achievement of control students is normally distributed and the treatment shifts the mean by $d$ standard deviations, a treated student of average ability would be at the $\Phi(d)$ quantile of the control distribution. For $d=2$ this is $\Phi(2)\approx0.977$, and for $d=0.8$ it is $\Phi(0.8)\approx0.788$. The program uses the lower value for planning, since it is better supported, and treats the higher as the result of conditions that may not be reproducible at scale. Whatever value is used, the benefit is reported together with its measurement protocol, namely retention and transfer assessed without assistance after a stated delay.
